\documentclass[10pt,a4paper]{article}
\usepackage{amsmath,amssymb,amsthm,graphicx}
\usepackage{xfrac}
\newenvironment{enumeratei}
{\bgroup\def\theenumi{\roman{enumi}}\def\theenumii{\arabic{enumii}}\begin{enumerate}}
{\end{enumerate}\egroup}

\hyphenation{addi-tion-al belong belongs denote exchan-ge exchan-ged exchan-ges}

\newcommand{\psfrac}[2]{\sfrac{(#1)}{#2}}

\newcommand\mto{\mathchoice{\longmapsto}{\mapsto}{\mapsto}{\mapsto}}
\let\oldbigsqcup\bigsqcup
\renewcommand{\bigsqcup}{\mathop{\textstyle\oldbigsqcup}\displaylimits}

\let\epsilon\varepsilon

\DeclareMathOperator{\Ess}{Ess}
\DeclareMathOperator{\Leb}{Leb}

\usepackage{tikz}
\usetikzlibrary{plotmarks}

\theoremstyle{plain}
\newtheorem{theorem}{Theorem}[section]
\newtheorem{lemma}[theorem]{Lemma}
\newtheorem{proposition}[theorem]{Proposition}
\newtheorem{corollary}[theorem]{Corollary}
\newtheorem{claim}{Claim}

\theoremstyle{definition}
\newtheorem{remark}[theorem]{Remark}
\newtheorem{example}[theorem]{Example}

\newcommand{\A}{\mathcal{A}}
\newcommand{\B}{\mathcal{B}}
\newcommand{\I}{\mathcal{I}}

\newcommand{\Z}{\mathbb{Z}}
\newcommand{\N}{\mathbb{N}}
\newcommand{\R}{\mathbb{R}}
\title{The first backward return of a centre to a symmetric subinterval lands on a centre or on an endpoint}
\author{Przemyslaw Berk, Frank Trujillo}
\date{doi:10.5802/jep.302; CC BY4.0}
\begin{document}
\maketitle
\noindent\textit{Editorial note.} The counted statement is the author's labelled claim environment at source0.tex lines 761-775; its proof environment opens at line 777, the line after the blank separator, so the pairing is by adjacency and the retained proof range 777-800 is the whole environment.\par\medskip
\noindent\textit{Editorial note.} Five statements are retained as context because the statement or the proof refers to them: eq:generalized\_conjugacy (source lines 469-472), eq:inverse\_iterates (source lines 533-537), eq:delta\_1/2 (source lines 539-544), lem: local\_reflection (source lines 596-602) and lem:symmetric\_connection (source lines 652-670).\par\medskip
\noindent\textit{Editorial note.} This source is new to the material-authoring work and its preamble is built from the author's own macro and theorem block (source lines 8-45) rather than from a prepared adapter.\par\medskip
\noindent\textit{Editorial note.} This source has no label-number dump and the retained counter values are left empty, so the printed numbering of the card is the wrapper's default; the identity of the counted item is its author label.\par\medskip
\noindent\textit{Editorial note.} r002 note: r001 failed with undefined references to sc:induced\_IETs and lem: alphabetafetrconnections, both used inside the counted proof at source line 799; r002 adds exactly those two environments as further context and changes nothing else.\par\medskip
\noindent\textit{Editorial note.} r003 note: the r002 context lem: alphabetafetrconnections refers to eq: paramofJ, so r003 adds exactly that numbered description of J as a further context and changes nothing else.\par\medskip
\noindent\textit{Editorial scope.} One original result of the pinned publisher article, counted once: the dichotomy for the first backward return of a centre to a symmetric subinterval (author label cl:backward\_centers).\par\medskip
\noindent\textit{Allowed standard prerequisites (editorial).}\begin{itemize}
\item First backward return time to a subinterval. b\_J(c) is the first positive time the backward orbit of the point c enters the subinterval J.. It is the exponent governing the alternative in the conclusion.
\item Induced transformation and its centres. The induced transformation on J exchanges its continuity intervals, and the centres of those intervals are the points fixed by the conjugate of the induced transformation.. They are the possible values of the projected point in the first alternative.
\item Connections. A connection is a finite orbit fragment joining two discontinuities.. It is what the second alternative forces for the centre.
\end{itemize}\medskip
\begin{equation}\label{eq: paramofJ}
\begin{gathered}
\text{$a_J=T^{m_0}(\partial I_{\alpha_J})$ and $b_J=T^{n_0}(\partial I_{\beta_J})$, \quad for some $\alpha_J,\beta_J\in\mathcal A$ and $m_0,n_0\in\Z$,} \\
T^{m}(\partial I_{\alpha_J})\notin J, \qquad \text{for any } m \in \{0, \ldots, m_0\} \text{ with } m \neq m_0 \\
T^n(\partial I_{\beta_J})\notin J \qquad \text{for any } n \in \{0, \ldots, n_0\} \text{ with } n \neq n_0,
\end{gathered}
\end{equation}

\subsection{Induced IETs}\label{sc:induced_IETs}

\begin{lemma}\label{lem: alphabetafetrconnections}
Let $T$ be an ergodic interval exchange transformation of $d=\#\A$ intervals and let $J$ be a subinterval of the form \eqref{eq: paramofJ}. Then $T_J$ can be considered as an interval exchange of $d_J$ intervals, where
\[
d_J:= d - \#\{\alpha\in\A\setminus \{\pi_0^{-1}(1)\} \mid m_{J,\alpha}\ge M(\alpha)\}.
\]
In particular, if $J$ does not contain any point from any connection, then $d-d_J$ is equal to the number of non-trivial irreducible connections of $T$.
\end{lemma}

\begin{multline}
\label{eq:generalized_conjugacy}
\mathcal{I}_I\circ T^n(x)=T^{-n}\circ\mathcal{I}_I(x) \quad \text{if } x \neq T^{-i}(\partial I_\alpha) \text{ for any } \alpha \in \A \text{ and}\\ \text{any $i$ such that } 0 \leq i < n.
\end{multline}

\begin{equation}
\label{eq:inverse_iterates}
T^{-m}(c_\alpha) =
{T^{-1} \circ \I_I} \big( T^{m - \delta_{1/2}(\alpha)}(c_\alpha) \big),
\end{equation}

\begin{equation}
\label{eq:delta_1/2}
\delta_{1/2}(\beta) = \left\{\begin{array}{ll} 1 & \text{ if } \quad \beta = \sfrac{1}{2}, \\
0 & \text{ otherwise}.
\end{array} \right.
\end{equation}

\begin{lemma}\label{lem: local_reflection}
Let $T : I \to I$ be a symmetric IET and $J \subseteq I$ be an $\alpha$-symmetric interval, for some $\alpha \in \A$. Then $\mathcal{I}_I\circ T(\mathring{J}) = \mathring{J}$. Moreover,
\[
\I_J = \mathcal{I}_I\circ T|_{\mathring{J}},
\]
that is, $\mathcal{I}_I\circ T|_{\mathring{J}}$ is the symmetric reflection on $J$.
\end{lemma}

\begin{lemma}
\label{lem:symmetric_connection}
Let $T: I \to I$ be a symmetric IET and $\beta \in \A \cup \{\sfrac{1}{2}\}$. Assume there exists $m \in \Z$ and $\alpha \in \A$ such that $T^m(c_\beta) = \partial I_\alpha$ and $T^k(c_\beta) \notin \{\partial I_\gamma\}_{\gamma \in \A}$, for any $-|m| < k < |m|$.

Then {if $m\ge 0$}, we~have $\alpha \neq \pi_0^{-1}(1)$ and\vspace*{-3pt}
\[
T^{-m - \delta_{1/2}(\beta)}(c_\beta) = \partial I_{\hat \alpha},
\]
where $\delta_{1/2}$ is given by \eqref{eq:delta_1/2} and $\pi_0( \hat \alpha) = \pi_0(\alpha) - 1$.

{If on the other hand $m<0$, then $\alpha\neq \pi_1^{-1}(1)$ and\vspace*{-3pt}
\[
T^{-m - \delta_{1/2}(\beta)}(c_\beta) = \partial I_{\hat \alpha},
\]
with $\pi_0( \hat \alpha) = \pi_0(\alpha) + 1$.
}

In particular, $c_\beta$ lies inside a non-trivial connection.
\end{lemma}

\begin{claim}
\label{cl:backward_centers}
Let $\sigma \in \A \cup \{\sfrac{1}{2}\}$ with $\sigma \neq \beta$ and $\ell := b_J(c_\sigma) \geq 1$ be the first backwards return time of $c_\sigma$ to $J$, that is, such that $p_J(c_\sigma) = T^{-\ell}(c_\sigma)$. Then, the following dichotomy holds.
\begin{itemize}
\item either $\{c_\sigma, T(c_\sigma), \ldots, T^{\ell - \delta_{1/2}(\sigma)}(c_\sigma)\} \cap \{\partial I_\delta\}_{\delta \in \A} = \emptyset$ and
\[
p_J(c_\sigma) = \hat c_\gamma^J,\qquad T_J(\hat c_\gamma^J) = T^{\ell - \delta_{1/2}(\sigma)}(c_\sigma) = T^{2\ell - \delta_{1/2}(\sigma)}(\hat c_\gamma^J), \quad \text{for some } \gamma \in \mathcal C,
\]
\item or $\{c_\sigma, T(c_\sigma), \ldots, T^{\ell - \delta_{1/2}(\sigma))}(c_\sigma)\} \cap \{\partial I_\delta\}_{\delta \in \A} \neq \emptyset$ and
\[
p_J(c_\sigma) = \partial I_\gamma^J,\quad\text{for some } \gamma \in \A_J,
\]
\end{itemize}
where $\delta_{1/2}$ is given by \eqref{eq:delta_1/2}. Moreover, in the latter case, $c_\sigma$ lies in a non-trivial connection of $T$.
\end{claim}

\begin{proof}
First, let us assume that $\{c_\sigma, T(c_\sigma), \ldots, T^{\ell - \delta_{1/2}(\sigma)}(c_\sigma)\} \cap \{\partial I_\delta\}_{\delta \in \A} = \emptyset$.

By \eqref{eq:inverse_iterates} and Lemma \ref{lem: local_reflection},
$T^{-k}(c_\sigma) \in J$ if and only if $T^{k - \delta_{1/2}}(c_\sigma) \in J$,
for any $k$ such that $0 \leq k \leq \ell$. Hence, the first visit time of $c_\sigma$ to $J$
via $T$ is $\ell - \delta_{1/2}(\sigma)$, which implies
\[
T_J(p_J(c_\sigma)) = T^{\ell - \delta_{1/2}(\sigma)}(c_\sigma).
\]
Moreover, $p_J(c_\sigma) = T^{-\ell}(c_\sigma)$ is a fixed point of $\I_J \circ T_J$. Indeed, by~\eqref{eq:generalized_conjugacy} and Lemma~\ref{lem: local_reflection},
\begin{align*}
\I_J \circ T_J(p_J(c_\sigma)) &= \I \circ T \big( T^{\ell -
\delta_{1/2}(\sigma)}(c_\sigma) \big) \\
&= T^{-\ell + \delta_{1/2}(\sigma)}\circ
T^{-1} \circ \I (c_\sigma) = T^{-\ell + \delta_{1/2}(\sigma)} (c_\sigma) =
p_J(c_\sigma).
\end{align*}
Therefore, since $p_J(c_\sigma)$ is a fixed point of $\I_J \circ T_J$ and $T_J$ exchanges its continuity intervals symmetrically, $p_J(c_\sigma) = \hat c_\gamma^J$, for some $\gamma \in \mathcal C$.

Now, let us assume that $\{c_\sigma, T(c_\sigma), \ldots, T^{\ell - \delta_{1/2}(\sigma)}(c_\sigma)\} \cap \{\partial I_\delta\}_{\delta \in \A} \neq \emptyset$.

Let $k$ with $0 \leq k \leq \ell - \delta_{1/2}(\sigma)$ be the minimum such that $T^k(c_\sigma) = \partial I _\delta$ for some $\delta \in \A$. By Lemma \ref{lem:symmetric_connection}, $T^{-k - \delta_{1/2}(\sigma)}(c_\sigma) = \partial I_{\hat \delta}$, where $\pi_0(\hat \delta) = \pi_0(\delta) - 1$. Since $-k - \delta_{1/2}(\sigma) \geq -\ell$, it~follows that $p_J(\partial I_{\hat \delta}) = p_J(c_\sigma)$ which, by~definition of $T_J$ (see Section \ref{sc:induced_IETs} and Lemma \ref{lem: alphabetafetrconnections}), coincides with $\partial I_\gamma^J$, for some $\gamma \in \A_J$. Notice that, in~this case, $c_\sigma$ lies in a non-trivial connection of $T$.
\end{proof}

\end{document}
